\chapter{Úvod}
\label{1-uvod}

V~posledních desetiletích došlo k~výraznému rozvoji metod kosmické geodézie. Jednou z~těchto metod je metoda \zkratka{DORIS}, která je založena na měření Dopplerova posunu signálu mezi pozemními stanicemi, na nichž jsou umístěny vysílače, a družicemi, na nichž jsou umístěny přijímače. Výsledky této metody jsou využívány například pro určování přesných drah družic, určování polohy pozemních stanic, stanovení parametrů orientace Země nebo pro realizaci a zpřesňování mezinárodního terestrického referenčního rámce (\zk{ITRF})~\cite{ids_what_is_doris}. Koordinaci zpracování a distribuci produktů systému DORIS zajišťuje služba \zkratka{IDS}, která poskytuje jednotlivé produkty, tedy výsledky zpracování analytických center. Analytická centra \zkratka{IDS} zpracovávají data nezávisle různými metodickými postupy, a proto je nutné jejich výsledky pravidelně analyzovat a vzájemně porovnávat. Kvalita těchto výstupů má přímý vliv na přesnost určení souřadnic stanic i drah družic, a tím i na kvalitu globálních geodetických produktů.

Cílem této diplomové práce je analyzovat výsledky zpracování metody \zkratka{DORIS} s~důrazem na výstupy analytického centra na Geodetické observatoři Pecný (\zk{GOP}), která je součástí Výzkumného ústavu geodetického, topografického a kartogra\-fického, v.~v.~i. První část práce je věnována analýze časových řad souřadnic stanic vysílačů systému \zkratka{DORIS} za období 2015--2025. Analyzovány jsou stanice, u~nichž bylo v~interním přehledu \zkratka{IDS} indikováno podezřelé chování. Cílem je toto chování potvrdit nebo vyvrátit a případně se pokusit nalézt jeho příčinu. V~rámci analýzy jsou identifikovány dlouhodobé trendy, diskontinuity a periodické složky časových řad. Druhá část práce se zabývá analýzou přesných drah družic vybavených přijímačem systému \zkratka{DORIS}. Porovnávány jsou dráhy z~řešení analytického centra na Geodetické observatoři Pecný (\zk{GOP}) s~řešením francouzského analytického centra (\zk{SSA}) za prvních 30 dní roku 2024. Pro porovnání je využita Hermitova interpolace. Výsledky jsou ověřeny nezávislým porovnáním pomocí dynamické parametrizace dráhy v~\zkratka{BSW}. Cílem je posoudit shodu výsledků obou analytických center.

Výsledkem práce je zhodnocení výstupů metody \zkratka{DORIS} se zaměřením na řešení analytického centra na Geodetické observatoři Pecný. Pro účely obou analýz byla vytvořena knihovna v~jazyce Python, která umožňuje jednotné zpracování dat systému \zkratka{DORIS}, jejich analýzu a~generování výstupů. Kompletní výsledky analýz jednotlivých stanic a~družic jsou obsaženy v~reportech v~přílohách práce.